% =====================================================================
% preamble.tex
% Continuation Under Constraint: Distinctions, Substrates, and Repair
% Across Transformations
%
% Compilation: XeLaTeX is the tested, working engine for this
% skeleton. LuaLaTeX is the intended long-run primary stack (per
% established practice across the rest of the corpus), but requires
% the luaotfload package, which was found to be absent from this
% drafting environment's texlive install; fontspec fails at font-load
% time under lualatex here as a direct result. This is exactly the
% documented fallback scenario -- "XeLaTeX as fallback for
% font-loading issues" -- rather than an exception to it. On a
% machine with a complete texlive install (specifically the
% texlive-luatex collection, which provides luaotfload), switch back
% to lualatex by changing only the engine invocation; no preamble
% changes are required, since fontspec is used identically under both.
% =====================================================================

\usepackage{fontspec}
\setmainfont{TeX Gyre Pagella}
\usepackage[full]{textcomp}

% NOTE ON MATH FONT: unicode-math + a dedicated Pagella math font
% would be the ideal pairing, but the unicode-math stack requires
% lualatex-math.sty, which is not present in this environment's
% texlive install (and CTAN is outside the allowed network domains
% for fetching it here). mathpazo was tried as a fallback but its
% Type1/T1-encoded text fonts clash with fontspec's TU encoding under
% lualatex, so this skeleton falls back further to plain amsmath with
% default (Computer Modern) math symbols for now -- compiles cleanly
% under all three engines with zero encoding conflicts. Once drafting
% moves to an environment with a complete texlive install, swap in
% unicode-math + a Pagella-family math font (or newpxmath, which is
% fontspec-safe) for visual consistency with the Pagella body text;
% no chapter file needs to change for that swap.
\usepackage{amsmath}
\usepackage{amsthm}
\usepackage{amssymb}
\usepackage{mathtools}
\usepackage{microtype}
\usepackage[hidelinks]{hyperref}
\usepackage{cleveref}
\usepackage{enumitem}
\usepackage{geometry}
\geometry{a5paper,margin=2.2cm} % scholarly-monograph trim size; adjust for print run
\usepackage{titlesec}
\usepackage{fancyhdr}
\usepackage{makeidx}
\makeindex
\usepackage{booktabs}
\usepackage{longtable}

% ---------------------------------------------------------------------
% Numbering: theorem-like environments numbered within chapter, so
% cross-references remain stable as chapters are drafted independently
% and reordered during editing. This matches the practice used across
% the rest of the corpus (CPR, Semantic Infrastructure, etc.).
% ---------------------------------------------------------------------
\numberwithin{equation}{chapter}

\theoremstyle{plain}
\newtheorem{theorem}{Theorem}[chapter]
\newtheorem{proposition}[theorem]{Proposition}
\newtheorem{corollary}[theorem]{Corollary}
\newtheorem{lemma}[theorem]{Lemma}

\theoremstyle{definition}
\newtheorem{definition}[theorem]{Definition}
\newtheorem{construction}[theorem]{Construction}
\newtheorem{example}[theorem]{Example}

\theoremstyle{remark}
\newtheorem{remark}[theorem]{Remark}
\newtheorem*{proofsketch}{Proof Sketch}

% ---------------------------------------------------------------------
% Non-standard environments enforcing the monograph's epistemic-
% labeling discipline: a claim's status (proved theorem, corollary,
% proposition with a complete proof, or conjecture/open question
% awaiting fuller treatment) must be visually and semantically
% distinct. This is a structural commitment of the whole project,
% not a cosmetic choice -- see the Note on Epistemic Status in the
% front matter.
% ---------------------------------------------------------------------
\newtheorem{conjecture}[theorem]{Conjecture}
\newtheorem{openquestion}[theorem]{Open Question}
\newtheorem{principle}[theorem]{Principle} % informal, pre-formal statements (e.g. Ch. 1's thesis)

% Visual flag for imported results, so the reader can see at a glance
% which theorems originate in prior work (Conservation of Ambiguity,
% Observability Is Restorability, etc.) versus this monograph.
\newcommand{\imported}[1]{\textit{(Imported from #1.)}}
\newcommand{\newresult}{\textit{(New to this monograph.)}}
\newcommand{\transcribed}[1]{\textit{(Transcribed from #1 by the author;
not independently verified against primary source text.)}}
\newcommand{\secondary}[1]{\textit{(Reported from a prose summary of #1
by the author; neither transcribed formulas nor primary source text
independently available.)}}

% ---------------------------------------------------------------------
% Notation macros (populated as Part I is drafted; consolidated
% reference lives in Appendix B). Declared here so every chapter file
% can use them consistently without re-defining locally.
% ---------------------------------------------------------------------
\newcommand{\repcost}[1]{C_R(#1)}                 % repair cost
\newcommand{\amb}[1]{A(#1)}                        % operational ambiguity
\newcommand{\interf}[2]{I(#1,#2)}                  % repair interference
\newcommand{\discard}[1]{D(#1)}                    % discarded distinction set
\newcommand{\admis}[2]{\mathcal{A}_{#1}(#2)}       % admissible continuation set, task then substrate
\newcommand{\fidelity}[2]{F(#1,#2)}                % fidelity of rendering r for task tau
\newcommand{\restdomain}{\mathrm{Rest}(\mathcal{R})} % repair-stable domain, imported notation

\newcommand{\booktitle}{Continuation Under Constraint}
\newcommand{\booksubtitle}{Distinctions, Substrates, and Repair Across Transformations}

\pagestyle{fancy}
\fancyhf{}
\fancyhead[LE]{\textsc{\leftmark}}
\fancyhead[RO]{\textsc{\rightmark}}
\fancyfoot[C]{\thepage}
